\section{模型假设}
根据数据粒度、设备参数和购电规则，作如下假设与约定。
\begin{enumerate}
  \item 每个功率样本代表其标签之前十分钟区间的平均功率，区间电量按平均功率与区间时长之积计算，电价取对应样本值。
  \item 充放电量定义在微网母线侧，储电量定义在电池内部。题面90\%效率解释为充电效率和放电效率各为0.9，往返效率为0.81；5\,000 kW上限也施加于母线侧。
  \item 同一十分钟区间内储能只允许充电、放电或待机。效率取常数，除充放电转换损耗外，不另计自放电和线路损耗。
  \item 光伏可用于供负载、充电或弃光，不向外网售电。目标中只计题目规定的购电费用，未用合同电量仍付费，不另计电池寿命损耗费用。
  \item 问题一将附件1预测视为确定性输入，取日初储电量为6\,000 kWh，并要求日末恢复至该值。
  \item 问题四的计划仅使用当时已发布的预报与已结束区间的实际数据，预测在两次发布之间保持不变；实际供需用于执行保护和事后结算。假设本地控制器能在区间平均功率尺度下及时减少原定动作，实际储电量连续传递，不每日重置。
  \item 问题四以近期已实现联合误差近似规划期不确定性。三个情景共享全窗口合同与储能动作，供需分配可随情景变化；有限情景及固定权重只用于规划评价，不视为对所有实际扰动的覆盖保证。
\end{enumerate}
其中，效率解释决定储能递推及损耗大小；区间平均功率和及时响应假设限定了模型描述设备运行的时间尺度。问题四的回放初态与年度末态条件在相应章节另行说明。
